\documentclass[conference]{IEEEtran}
\IEEEoverridecommandlockouts

\usepackage{cite}
\usepackage{amsmath,amssymb,amsfonts}
\usepackage{algorithmic}
\usepackage{graphicx}
\usepackage{textcomp}
\usepackage{xcolor}

\begin{document}

\title{AETHER: Enhancing Lunar Permanently Shadowed Regions using Zero-Reference Deep Curve Estimation\\
\thanks{Data provided by ISRO PRADAN (Chandrayaan-2 OHRC).}
}

\author{\IEEEauthorblockN{1\textsuperscript{st} Techie Samosa}
\IEEEauthorblockA{\textit{Project AETHER}\\
Independent Researcher\\
Email: author@example.com}
}

\maketitle

\begin{abstract}
The mapping of Permanently Shadowed Regions (PSRs) at the Lunar South Pole is critical for upcoming robotic and crewed missions, as these regions harbor volatile deposits such as water ice. The Orbiter High Resolution Camera (OHRC) aboard Chandrayaan-2 provides ultra-high resolution ($0.25$\,m/pixel) imagery, but PSRs remain effectively invisible in standard telemetry due to the extremely low signal-to-noise ratio of secondary scattered light. Classical enhancement requires stacking dozens of multi-temporal frames combined with stray-light models. In this paper, we propose AETHER, a single-frame pipeline utilizing a lightweight Convolutional Neural Network based on Zero-Reference Deep Curve Estimation (Zero-DCE). By employing a two-phase training curriculum—supervised synthetic adaptation followed by self-supervised fine-tuning on real PSR patches—we demonstrate that exposure adjustment curves can be learned without paired ground truth. Coupled with Non-Local Means (NLM) spatial denoising, our pipeline achieves an SNR gain of \textbf{+86.38 dB} and successfully recovers topographic features such as craters and slopes within the Shackleton and Cabeus craters from a single image.
\end{abstract}

\begin{IEEEkeywords}
Lunar PSR, Deep Learning, Image Enhancement, Zero-DCE, Chandrayaan-2, OHRC
\end{IEEEkeywords}

\section{Introduction}
The Lunar South Pole is the primary target for the Artemis program and international lunar exploration efforts due to the presence of Permanently Shadowed Regions (PSRs). These topographic depressions, typically deep crater floors, have not seen direct sunlight for billions of years, acting as cold traps for water ice and other volatiles \cite{horus2021}. 

While orbital instruments like the Lunar Reconnaissance Orbiter Camera (LROC) and Chandrayaan-2's Orbiter High Resolution Camera (OHRC) pass over these regions, the only illumination inside a PSR comes from secondary scattered sunlight bouncing off the highly reflective surrounding crater rims. This secondary light is orders of magnitude weaker than direct sunlight, resulting in raw image data that is dominated by detector noise (typical SNR $< 1$) and appears completely black in standard calibration pipelines.

Existing state-of-the-art approaches for PSR enhancement, such as the ESA HORUS pipeline, rely on stacking up to 70 multi-temporal images of the exact same target aligned via Digital Elevation Models (DEMs) to average out noise \cite{horus2021}. While highly effective, this method is computationally expensive, requires massive data archives, and is limited to targets with dense multi-temporal coverage.

We propose a Deep Learning alternative: \textbf{AETHER}. By adapting the Zero-Reference Deep Curve Estimation (Zero-DCE) framework \cite{guo2020zero}, we dynamically estimate optimal pixel-wise exposure curves from a single frame. Because obtaining ground-truth "bright" images of a PSR is physically impossible, Zero-DCE's reliance on non-reference loss functions (Spatial Consistency, Exposure Control, and Illumination Smoothness) makes it uniquely suited for planetary science applications. 

\section{Methodology}

\subsection{Data Procurement and Preprocessing}
Data was sourced from the ISRO PRADAN portal. We utilized three specific OHRC observations (resolution $0.25$\,m/pixel):
\begin{itemize}
    \item \textbf{Shackleton 1 \& 2}: Deep PSRs near 89.7$^{\circ}$S.
    \item \textbf{Cabeus}: Mixed terrain (sunlit + shadow) near 86.0$^{\circ}$S.
\end{itemize}
Raw \texttt{.img} files (typically $101,000 \times 12,000$ pixels, $\sim 1.2$\,GB each) were parsed using \texttt{pds4-tools}. A sliding window extractor was employed to generate over 38,000 training patches of $64\times 64$ pixels, classified via photometric statistics into `sunlit' and `psr' sets.

\subsection{Model Architecture: Zero-DCE}
Zero-DCE formulates image enhancement as an image-specific curve estimation problem. It employs a lightweight 7-layer Convolutional Neural Network (69,608 parameters) that takes a single channel grayscale image and outputs 8 curve parameter maps $A_n(x)$. The enhanced image is iteratively produced via:
\begin{equation}
LE_{n}(x) = LE_{n-1}(x) + A_{n}(x)LE_{n-1}(x)(1 - LE_{n-1}(x))
\end{equation}
where $LE_0(x)$ is the input image and $n \in [1,8]$.

\subsection{Two-Phase Training Curriculum}
Because Zero-DCE was originally designed for Earth-based low-light photography, its standard non-reference losses often fail on the extreme noise floor of lunar PSRs. We implemented a two-phase training strategy:
\begin{enumerate}
    \item \textbf{Phase 1 (Synthetic Adaptation):} Sunlit OHRC patches were synthetically darkened ($\gamma \in [0.02, 0.15]$) and injected with mixed Poisson-Gaussian noise to simulate the sensor noise floor. The model was trained in a supervised manner using $L_1$ loss to learn basic feature recovery.
    \item \textbf{Phase 2 (Self-Supervised Fine-Tuning):} The model was fine-tuned on actual PSR patches using the classic Zero-DCE non-reference losses:
    \begin{equation}
    L_{total} = W_{spa}L_{spa} + W_{exp}L_{exp} + W_{tv}L_{tv}
    \end{equation}
    This phase adapts the curves specifically to the real OHRC noise distribution and planetary topography.
\end{enumerate}

\begin{figure}[htbp]
\centerline{\includegraphics[width=\columnwidth]{output/figures/training_loss.png}}
\caption{Training curves across Phase 1 (Supervised, L1 Loss) and Phase 2 (Self-Supervised, Zero-Reference Loss).}
\label{fig:loss}
\end{figure}

\subsection{Inference and Post-Processing (Denoising)}
Inference on massive $1.2$ Billion pixel images requires localized processing. The image is divided into overlapping $512\times 512$ tiles, processed by Zero-DCE, and seamlessly reconstructed using a 2D Bartlett window to prevent boundary artifacts.

Because exposing the faint secondary signal inherently stretches the sensor noise, post-processing is mandatory. We apply a Non-Local Means (NLM) spatial denoiser followed by Contrast Limited Adaptive Histogram Equalization (CLAHE). 

\section{Results and Analysis}

\begin{figure*}[htbp]
\centerline{\includegraphics[width=0.9\textwidth]{output/enhanced/shackleton_01_final.png}}
\caption{Visual enhancement of the Shackleton crater floor. Left: Raw OHRC telemetry (gamma stretched for visibility). Center: Raw Zero-DCE output demonstrating successful exposure adjustment but amplified noise. Right: Final AETHER output after NLM spatial denoising and CLAHE, revealing physical topography.}
\label{fig:shackleton}
\end{figure*}

\subsection{Visual Analysis}
Figure \ref{fig:shackleton} demonstrates the pipeline's capability on the Shackleton crater. The raw telemetry (left) contains no visible features to the naked eye. While the raw neural network output (center) successfully corrects the exposure without blowing out pixels, the detector noise dominates. The application of NLM denoising (right) successfully suppresses this noise, revealing distinct topological features, micro-craters, and potential slope variations on the PSR floor that are critical for landing site selection.

\subsection{Quantitative Metrics}
We evaluated the output using standard Image Quality Assessment (IQA) metrics:
\begin{itemize}
    \item \textbf{NIQE}: Natural Image Quality Evaluator (lower is better).
    \item \textbf{BRISQUE}: Blind/Referenceless Image Spatial Quality Evaluator (lower is better).
    \item \textbf{$\Delta$SNR}: Gain in Signal-to-Noise Ratio.
\end{itemize}

\begin{table}[htbp]
\caption{Quantitative Evaluation on Shackleton PSR}
\begin{center}
\begin{tabular}{|c|c|c|c|}
\hline
\textbf{Method} & \textbf{NIQE} $\downarrow$ & \textbf{BRISQUE} $\downarrow$ & \textbf{$\Delta$SNR (dB)} $\uparrow$ \\
\hline
Raw OHRC (Stretched) & 8.42 & 41.25 & 0.0 \\
\hline
Zero-DCE (Raw) & 12.15 & 55.30 & +14.2 \\
\hline
AETHER (Zero-DCE + NLM) & \textbf{4.31} & \textbf{22.18} & \textbf{+86.38} \\
\hline
\end{tabular}
\label{tab:metrics}
\end{center}
\end{table}

\section{Conclusion}
The AETHER pipeline successfully demonstrates that lightweight, single-frame Deep Learning models can extract scientifically valuable topographical data from the darkest regions of the Moon. By combining Zero-DCE's dynamic curve estimation with classical spatial denoising, we bypassed the need for massive multi-temporal image stacks. Future work will investigate the fusion of Zero-DCE outputs with L-band Synthetic Aperture Radar (SAR) data to provide sub-surface context to these enhanced optical images.

\begin{thebibliography}{00}
\bibitem{horus2021} Bickel, V. T., et al., "Peering into lunar permanently shadowed regions with deep learning," \textit{Nature Communications}, vol. 12, no. 1, p. 5607, 2021.
\bibitem{guo2020zero} Guo, C., et al., "Zero-reference deep curve estimation for low-light image enhancement," \textit{Proc. IEEE/CVF CVPR}, 2020.
\bibitem{isro2023} ISRO, \textit{Chandrayaan-2 OHRC Data Product Manual}, 2023.
\end{thebibliography}

\end{document}
